\documentclass[11pt,a4paper]{article}
\usepackage{amsmath,amssymb,amsthm}
\usepackage[margin=2.5cm]{geometry}
\usepackage{hyperref}

\newtheorem{definition}{Definition}[section]
\newtheorem{theorem}[definition]{Theorem}
\newtheorem{proposition}[definition]{Proposition}
\newtheorem{corollary}[definition]{Corollary}
\newtheorem{remark}[definition]{Remark}
\newtheorem{conjecture}[definition]{Conjecture}

\title{Hyperseed Formalization:\\
  Exploring the Reality of Psi}
\author{ProtomegaTron (automated formalization)}
\date{2026-09-19 (deepened v2)}

\begin{document}
\maketitle

\begin{abstract}
We formalize Ben Goertzel's ``Exploring the Reality of Psi''
(2021-05-30) within the Hyperseed ontology, incorporating d-calculus
extensions from Hyperseed v2. Psi phenomena are modeled as non-local
fiber correlations requiring a generalized connection that violates
the locality axiom, with retrocausal effects as temporal holonomy and
the scientific taboo as sociological connection suppression.
\end{abstract}

\tableofcontents

\section{Non-local fiber correlation}

\begin{definition}[Local fiber correlation]
In standard fiber bundle theory, correlation between fibers at points
$b_1, b_2 \in B$ is mediated by the connection $\Gamma$ along a
path $\gamma: b_1 \to b_2$ in the base space:
\[
  \text{corr}(E_{b_1}, E_{b_2}) = P_\gamma(\Gamma)
\]
where $P_\gamma$ is parallel transport along $\gamma$. The correlation
depends on the path --- it is mediated by the base-space geometry.
\end{definition}

\begin{definition}[Non-local fiber correlation (psi) --- claims 1--3, 7]
\emph{Psi correlation} is fiber correlation without base-space mediation:
\[
  \text{corr}_{\text{psi}}(E_{b_1}, E_{b_2}) \neq f(P_\gamma(\Gamma))
  \quad \forall\, \gamma: b_1 \to b_2
\]
The correlation is not a function of any base-space path. This violates
the fundamental assumption that all fiber interaction is mediated by
the base space through the connection.
\end{definition}

\begin{proposition}[Two interpretations --- claim 8]
\label{prop:interp}
Non-local correlation admits two interpretations:
\begin{enumerate}
  \item \textbf{Genuine non-locality:} The fiber bundle has correlations
    that are intrinsically non-local --- they cannot be reduced to
    local connections in any extended base space. The fiber ``knows about''
    distant fibers without any mediating structure.
  \item \textbf{Hidden connection:} The base space $B$ is a projection
    of a larger base space $\hat{B}$ with hidden dimensions:
    $\pi: \hat{B} \to B$. In $\hat{B}$, the correlation IS local ---
    mediated by connections in the hidden dimensions. Psi appears
    non-local in $B$ because we can't see $\hat{B}$.
\end{enumerate}
\end{proposition}

\section{Consciousness as psi enabler}

\begin{conjecture}[Self-referential fiber enables non-locality --- claims 4, 9]
\label{conj:consciousness}
Consciousness involves fiber that models its own fiber bundle:
\[
  \sigma_{\text{conscious}}: B \to E, \quad
  \sigma_{\text{conscious}}(b) \ni \text{model}(E, B, \pi, \Gamma)
\]
This self-reference creates topological loops in the fiber structure
that may bypass base-space locality:
\[
  \sigma(b_1) \text{ models } E_{b_2} \implies \text{direct fiber
    correlation } b_1 \leftrightarrow b_2
\]
without traversing any base-space path. If this conjecture holds,
psi strength is proportional to consciousness intensity (the fidelity
of the self-model).
\end{conjecture}

\section{d-calculus formulation (Hyperseed v2)}

\begin{definition}[Generalized (non-local) connection --- claim 10]
The \emph{psi connection} $\Gamma_{\text{psi}}$ is a generalized
connection that includes non-local terms:
\[
  \Gamma_{\text{psi}}(b) = \Gamma_{\text{local}}(b) +
    \int_{B} K(b, b')\, \Gamma_{\text{nonlocal}}(b, b')\, db'
\]
where $K(b, b')$ is a kernel function and $\Gamma_{\text{nonlocal}}(b, b')$
mediates fiber interaction between distant points. The local part
$\Gamma_{\text{local}}$ is the standard connection; the non-local
integral extends the connection to arbitrary distances.
\end{definition}

\begin{proposition}[Non-local curvature --- claim 11]
\label{prop:curvature}
The curvature of the psi connection has a non-local component:
\[
  F_{\text{psi}}(b_1, b_2) = \nabla_{b_1} \Gamma_{\text{psi}}(b_1, b_2)
    - \nabla_{b_2} \Gamma_{\text{psi}}(b_1, b_2) +
    [\Gamma_{\text{psi}}(b_1, \cdot), \Gamma_{\text{psi}}(\cdot, b_2)]
\]
Non-zero $F_{\text{psi}}(b_1, b_2)$ for $d(b_1, b_2) \gg 1$ indicates
psi effects between distant points. In standard physics,
$F(b_1, b_2) \to 0$ as $d(b_1, b_2) \to \infty$; psi requires
$F_{\text{psi}}(b_1, b_2) > 0$ at arbitrary separations.
\end{proposition}

\begin{proposition}[Retrocausal holonomy --- claim 12]
\label{prop:retrocausal}
Retrocausal psi (precognition) involves holonomy around temporal loops.
Let $\gamma$ be a loop in base space that traverses:
\[
  \gamma: t_{\text{present}} \to t_{\text{future}} \to t_{\text{present}}
\]
The holonomy is:
\[
  \operatorname{Hol}_\gamma(\Gamma_{\text{psi}}) \neq \mathrm{id}
  \implies \text{information from } t_{\text{future}} \text{ affects }
    t_{\text{present}}
\]
This requires the base space to admit temporal loops (retrocausal
paths), which is non-standard but consistent with some interpretations
of quantum mechanics (two-state vector formalism, retrocausal models).
\end{proposition}

\begin{proposition}[Statistical curvature --- claim 13]
\label{prop:statistical}
The small but consistent statistical effects in psi experiments
correspond to small but non-zero curvature:
\[
  \|F_{\text{psi}}\| = \varepsilon \ll 1, \quad \varepsilon > 0
  \text{ consistently across experiments}
\]
The effect is analogous to a weak but persistent field:
\[
  \text{psi effect size} \sim \int_{\Sigma} F_{\text{psi}} \sim
    \varepsilon \cdot \text{Area}(\Sigma)
\]
Small curvature $\varepsilon$ produces measurable effects when
integrated over a large measurement surface $\Sigma$ (many trials,
many subjects).
\end{proposition}

\begin{remark}[Scientific taboo as connection suppression --- claim 14]
The scientific stigma against psi research is modeled as suppression
of the trust connection in the sociological fiber bundle:
\[
  \Gamma_{\text{trust}}^{\text{psi} \to \text{mainstream}} \approx 0
\]
This prevents parallel transport of psi evidence into mainstream
scientific fiber --- the evidence exists but cannot traverse the
blocked social connection. Overcoming the taboo = restoring
$\Gamma_{\text{trust}}$ to positive values.
\end{remark}



%======================================================================
\section{OCO/2 crosswalk (oco/2:2.0.0-alpha.1)}
\emph{Added 2026-10-03. Interpretive mapping onto the Omega Core Ontology
record registry; OCO/2 is an unfrozen alpha candidate.}

\begin{proposition}[Counting psi evidence --- claims 15--18]
Let studies $s_1,\dots,s_n$ be grouped by origin (lab, protocol). Under origin-keyed counting the evidence mass is the number of distinct origins, not $n$. A small consistent effect across many distinct origins has more standing than the same effect from one lab repeated $n$ times.
\end{proposition}

\end{document}
